\documentclass[../main.tex]{subfiles}
\begin{document}

\section{3.3.}

Let $W$ be a random variable giving the number of heads minus the number
of tailes in three tosses of a coin. List the elements of the sample space S
for the three tosses of the coin and to each sample point assgin
a value $w$ of $W$.

\vspace{15em}

\section{3.8.}
Find the probability distribution of the random variable $W$ in \textbf{Exercise 3.3},
assuming that the coin is biased so that a head is twice as likely to occur as a tail.

\vspace{20em}

\section{3.23.}
Find the cumulative distribution function of the random variable
$W$ in \textbf{Exercise 3.8}. Using $F(w)$, find
\\\\
(a) $P(W>0)$;
\\\\
(b) $P(-1\leq W < 3)$.

\vspace{15em}

\section{3.11.}
A shipment of 7 televisions contains 2 defective ones. A hotel makes
a random purchase of 3 televisions. If $x$ is the number of defective ones
purchased by the hotel, find the probability distribution of $X$.
Express the results graphically as a probability histogram.

\vspace{15em}

\section{3.15}
Find the cumulative distribution function of the random variable $X$ representing 
the number of defectives in \textbf{Exercise 3.11}. Then using $F(x)$, find
\\\\
(a) $P(X=1)$;
\\\\
(b) $P(0<X\leq 2)$.

\vspace{15em}

\section{3.16.}
Contruct a graph of the cumulative distribution funciton of \textbf{Exercise 3.15}.

\vspace{23em}

\section{3.14.}
The waiting time, in hours, between successive speeders spotted by a radar unit is a
continuous random variable with a cumulative distribution function
$$
F(x)=
\begin{cases}
0 & x<0 \\
1-e^{-4x} & x\geq 0
\end{cases}
$$
Find the probability of waiting fewer than 12 minutes between 
successive speeders
\\\\
(a) using the cumulative distribution function of $X$.
\\\\
(b) using the probability density function of $X$. 

\vspace{23em}

\section{3.29.}
An important factor in solid missile fuel is the particle size distribution.
Significant problems occur if the particle sizes are too large. From production data 
in the past, it has been determined that the particle size \(in micrometers\) distribution 
characterized by
$$
f(x)=
\begin{cases}
3x^{-4}, & x>1, \\
0, & \text{elsewhere}.
\end{cases}
$$
\\\\
(a) Verify that this is a valid probability function.
\\\\
(b) Evaluate $F(x)$.
\\\\
(c) What is the probability that a random particle from
the manufactured fuel exceeds 4 micrometers?
\end{document}